% =========================================================
% Section: Boolean Algebra of Sets
% Chapter: Algebras of Sets — Volume I
% =========================================================

\begin{tcolorbox}[colback=gray!6, colframe=gray!40, arc=2pt,
  left=6pt, right=6pt, top=4pt, bottom=4pt,
  title={\small\textbf{Section Toolkit --- Quick Reference}},
  fonttitle=\small\bfseries]
\small
\begin{tabular}{l l l l}
\toprule
\textbf{Label} & \textbf{Statement} & \textbf{Proof method} & \textbf{Detail} \\
\midrule
P\ref{prop:comm-union}
  & $A\cup B = B\cup A$
  & Element chase & \hyperref[prf:comm-union]{$\downarrow$} \\
P\ref{prop:comm-intersection}
  & $A\cap B = B\cap A$
  & Element chase & \hyperref[prf:comm-intersection]{$\downarrow$} \\
P\ref{prop:assoc-union}
  & $A\cup(B\cup C) = (A\cup B)\cup C$
  & Element chase & \hyperref[prf:assoc-union]{$\downarrow$} \\
P\ref{prop:assoc-intersection}
  & $A\cap(B\cap C) = (A\cap B)\cap C$
  & Element chase & \hyperref[prf:assoc-intersection]{$\downarrow$} \\
P\ref{prop:dist-cap-over-cup}
  & $A\cap(B\cup C) = (A\cap B)\cup(A\cap C)$
  & Element chase & \hyperref[prf:dist-cap-over-cup]{$\downarrow$} \\
P\ref{prop:dist-cup-over-cap}
  & $A\cup(B\cap C) = (A\cup B)\cap(A\cup C)$
  & Element chase & \hyperref[prf:dist-cup-over-cap]{$\downarrow$} \\
P\ref{prop:demorgan-cup}
  & $(A\cup B)^c = A^c\cap B^c$
  & Element chase & \hyperref[prf:demorgan-cup]{$\downarrow$} \\
P\ref{prop:demorgan-cap}
  & $(A\cap B)^c = A^c\cup B^c$
  & Element chase & \hyperref[prf:demorgan-cap]{$\downarrow$} \\
P\ref{prop:identity-union}
  & $A\cup\emptyset = A$
  & Element chase & \hyperref[prf:identity-union]{$\downarrow$} \\
P\ref{prop:identity-intersection}
  & $A\cap X = A$
  & Element chase & \hyperref[prf:identity-intersection]{$\downarrow$} \\
P\ref{prop:complement-union}
  & $A\cup A^c = X$
  & Element chase & \hyperref[prf:complement-union]{$\downarrow$} \\
P\ref{prop:complement-intersection}
  & $A\cap A^c = \emptyset$
  & Element chase & \hyperref[prf:complement-intersection]{$\downarrow$} \\
P\ref{prop:idempotence-union}
  & $A\cup A = A$
  & Element chase & \hyperref[prf:idempotence-union]{$\downarrow$} \\
P\ref{prop:idempotence-intersection}
  & $A\cap A = A$
  & Element chase & \hyperref[prf:idempotence-intersection]{$\downarrow$} \\
P\ref{prop:double-complement}
  & $(A^c)^c = A$
  & Element chase & \hyperref[prf:double-complement]{$\downarrow$} \\
P\ref{prop:absorption-cup}
  & $A\cup(A\cap B) = A$
  & Element chase & \hyperref[prf:absorption-cup]{$\downarrow$} \\
P\ref{prop:absorption-cap}
  & $A\cap(A\cup B) = A$
  & Element chase & \hyperref[prf:absorption-cap]{$\downarrow$} \\
P\ref{prop:symmdiff-comm}
  & $A\triangle B = B\triangle A$
  & Unfold definition & \hyperref[prf:symmdiff-comm]{$\downarrow$} \\
P\ref{prop:symmdiff-assoc}
  & $A\triangle(B\triangle C) = (A\triangle B)\triangle C$
  & Element chase (hard) & \hyperref[prf:symmdiff-assoc]{$\downarrow$} \\
P\ref{prop:symmdiff-self}
  & $A\triangle A = \emptyset$
  & Unfold definition & \hyperref[prf:symmdiff-self]{$\downarrow$} \\
P\ref{prop:symmdiff-empty}
  & $A\triangle\emptyset = A$
  & Unfold definition & \hyperref[prf:symmdiff-empty]{$\downarrow$} \\
P\ref{prop:cap-dist-symmdiff}
  & $A\cap(B\triangle C)=(A\cap B)\triangle(A\cap C)$
  & Element chase & \hyperref[prf:cap-dist-symmdiff]{$\downarrow$} \\
P\ref{prop:powerset-boolean-ring}
  & $(\mathcal{P}(X),\triangle,\cap)$ is a Boolean ring
  & Verify ring axioms & \hyperref[prf:powerset-boolean-ring]{$\downarrow$} \\
\bottomrule
\end{tabular}
\end{tcolorbox}

% ---------------------------------------------------------
% Commutativity
% ---------------------------------------------------------

\begin{proposition}[\texorpdfstring{%
  \hyperref[prf:comm-union]{Commutativity of Union}}{Commutativity of Union}]
\label{prop:comm-union}
For all sets $A, B \subseteq X$: $A \cup B = B \cup A$.
\end{proposition}

\begin{remark*}[Fully quantified form]
$\forall A,B\subseteq X:\ A\cup B = B\cup A$.
\end{remark*}

\begin{remark*}[Flash]
\FlashQ{Prove commutativity of union from the element-chase definition.}
\FlashA{$x\in A\cup B \Leftrightarrow x\in A$ or $x\in B
  \Leftrightarrow x\in B$ or $x\in A \Leftrightarrow x\in B\cup A$.}
\FlashImplies{prop:assoc-union}
\end{remark*}

\begin{proposition}[\texorpdfstring{%
  \hyperref[prf:comm-intersection]{Commutativity of Intersection}}{%
  Commutativity of Intersection}]
\label{prop:comm-intersection}
For all sets $A, B \subseteq X$: $A \cap B = B \cap A$.
\end{proposition}

\begin{remark*}[Fully quantified form]
$\forall A,B\subseteq X:\ A\cap B = B\cap A$.
\end{remark*}

\begin{remark*}[Flash]
\FlashQ{Prove commutativity of intersection from the element-chase definition.}
\FlashA{$x\in A\cap B \Leftrightarrow x\in A$ and $x\in B
  \Leftrightarrow x\in B$ and $x\in A \Leftrightarrow x\in B\cap A$.}
\FlashImplies{prop:assoc-intersection}
\end{remark*}

% ---------------------------------------------------------
% Associativity
% ---------------------------------------------------------

\begin{proposition}[\texorpdfstring{%
  \hyperref[prf:assoc-union]{Associativity of Union}}{Associativity of Union}]
\label{prop:assoc-union}
For all sets $A, B, C \subseteq X$:
$A \cup (B \cup C) = (A \cup B) \cup C$.
\end{proposition}

\begin{remark*}[Fully quantified form]
$\forall A,B,C\subseteq X:\ A\cup(B\cup C) = (A\cup B)\cup C$.
\end{remark*}

\begin{remark*}[Flash]
\FlashQ{Outline the element-chase proof of associativity of union.}
\FlashA{$x\in A\cup(B\cup C)
  \Leftrightarrow x\in A$ or ($x\in B$ or $x\in C$)
  $\Leftrightarrow$ ($x\in A$ or $x\in B$) or $x\in C$
  $\Leftrightarrow x\in(A\cup B)\cup C$.
  The middle step is associativity of logical disjunction.}
\FlashImplies{prop:powerset-boolean-ring}
\end{remark*}

\begin{proposition}[\texorpdfstring{%
  \hyperref[prf:assoc-intersection]{Associativity of Intersection}}{%
  Associativity of Intersection}]
\label{prop:assoc-intersection}
For all sets $A, B, C \subseteq X$:
$A \cap (B \cap C) = (A \cap B) \cap C$.
\end{proposition}

\begin{remark*}[Fully quantified form]
$\forall A,B,C\subseteq X:\ A\cap(B\cap C) = (A\cap B)\cap C$.
\end{remark*}

\begin{remark*}[Flash]
\FlashQ{Outline the element-chase proof of associativity of intersection.}
\FlashA{$x\in A\cap(B\cap C)
  \Leftrightarrow x\in A$ and ($x\in B$ and $x\in C$)
  $\Leftrightarrow$ ($x\in A$ and $x\in B$) and $x\in C$
  $\Leftrightarrow x\in(A\cap B)\cap C$.
  The middle step is associativity of logical conjunction.}
\FlashImplies{prop:powerset-boolean-ring}
\end{remark*}

% ---------------------------------------------------------
% Distributivity
% ---------------------------------------------------------

\begin{proposition}[\texorpdfstring{%
  \hyperref[prf:dist-cap-over-cup]{%
    Distributivity of $\cap$ over $\cup$}}{%
  Distributivity of Intersection over Union}]
\label{prop:dist-cap-over-cup}
For all sets $A, B, C \subseteq X$:
$A \cap (B \cup C) = (A \cap B) \cup (A \cap C)$.
\end{proposition}

\begin{remark*}[Fully quantified form]
$\forall A,B,C\subseteq X:\
  A\cap(B\cup C) = (A\cap B)\cup(A\cap C)$.
\end{remark*}

\begin{remark*}[Flash]
\FlashQ{State and outline the proof of distributivity of $\cap$ over $\cup$.}
\FlashA{$A\cap(B\cup C) = (A\cap B)\cup(A\cap C)$.
  Element chase: $x$ is in the left side iff $x\in A$ and ($x\in B$ or $x\in C$),
  iff ($x\in A$ and $x\in B$) or ($x\in A$ and $x\in C$), iff $x$ is in the right.}
\FlashImplies{prop:powerset-boolean-ring}
\end{remark*}

\begin{proposition}[\texorpdfstring{%
  \hyperref[prf:dist-cup-over-cap]{%
    Distributivity of $\cup$ over $\cap$}}{%
  Distributivity of Union over Intersection}]
\label{prop:dist-cup-over-cap}
For all sets $A, B, C \subseteq X$:
$A \cup (B \cap C) = (A \cup B) \cap (A \cup C)$.
\end{proposition}

\begin{remark*}[Fully quantified form]
$\forall A,B,C\subseteq X:\
  A\cup(B\cap C) = (A\cup B)\cap(A\cup C)$.
\end{remark*}

\begin{remark*}[Flash]
\FlashQ{State and outline the proof of distributivity of $\cup$ over $\cap$.}
\FlashA{$A\cup(B\cap C) = (A\cup B)\cap(A\cup C)$.
  Element chase: $x$ is in the left side iff $x\in A$ or ($x\in B$ and $x\in C$),
  iff ($x\in A$ or $x\in B$) and ($x\in A$ or $x\in C$), iff $x$ is in the right.}
\FlashImplies{prop:powerset-boolean-ring}
\end{remark*}

% ---------------------------------------------------------
% De Morgan
% ---------------------------------------------------------

\begin{proposition}[\texorpdfstring{%
  \hyperref[prf:demorgan-cup]{De Morgan --- Union}}{De Morgan --- Union}]
\label{prop:demorgan-cup}
For all sets $A, B \subseteq X$: $(A \cup B)^c = A^c \cap B^c$.
\end{proposition}

\begin{remark*}[Fully quantified form]
$\forall A,B\subseteq X:\ (A\cup B)^c = A^c\cap B^c$.
\end{remark*}

\begin{remark*}[Flash]
\FlashQ{Prove De Morgan's law for union.}
\FlashA{$x\in(A\cup B)^c \Leftrightarrow x\notin A\cup B
  \Leftrightarrow \neg(x\in A\text{ or }x\in B)
  \Leftrightarrow x\notin A\text{ and }x\notin B
  \Leftrightarrow x\in A^c\cap B^c$.}
\FlashHint{The key step is De Morgan for logic: $\neg(P\vee Q)\Leftrightarrow\neg P\wedge\neg Q$.}
\FlashImplies{prop:sigma-closed-countable-intersection}
\end{remark*}

\begin{proposition}[\texorpdfstring{%
  \hyperref[prf:demorgan-cap]{De Morgan --- Intersection}}{%
  De Morgan --- Intersection}]
\label{prop:demorgan-cap}
For all sets $A, B \subseteq X$: $(A \cap B)^c = A^c \cup B^c$.
\end{proposition}

\begin{remark*}[Fully quantified form]
$\forall A,B\subseteq X:\ (A\cap B)^c = A^c\cup B^c$.
\end{remark*}

\begin{remark*}[Flash]
\FlashQ{Prove De Morgan's law for intersection.}
\FlashA{$x\in(A\cap B)^c \Leftrightarrow x\notin A\cap B
  \Leftrightarrow \neg(x\in A\text{ and }x\in B)
  \Leftrightarrow x\notin A\text{ or }x\notin B
  \Leftrightarrow x\in A^c\cup B^c$.}
\FlashHint{Key step: $\neg(P\wedge Q)\Leftrightarrow\neg P\vee\neg Q$.}
\FlashImplies{prop:sigma-closed-countable-intersection}
\end{remark*}

% ---------------------------------------------------------
% Identity, Complement, Idempotence
% ---------------------------------------------------------

\begin{proposition}[\texorpdfstring{%
  \hyperref[prf:identity-union]{Identity for Union}}{Identity for Union}]
\label{prop:identity-union}
For all $A \subseteq X$: $A \cup \emptyset = A$.
\end{proposition}

\begin{remark*}[Fully quantified form]
$\forall A\subseteq X:\ A\cup\emptyset = A$.
\end{remark*}

\begin{remark*}[Flash]
\FlashQ{Prove $A\cup\emptyset = A$.}
\FlashA{$x\in A\cup\emptyset \Leftrightarrow x\in A$ or $x\in\emptyset
  \Leftrightarrow x\in A$ or $\bot \Leftrightarrow x\in A$.}
\FlashImplies{prop:powerset-boolean-ring}
\end{remark*}

\begin{proposition}[\texorpdfstring{%
  \hyperref[prf:identity-intersection]{Identity for Intersection}}{%
  Identity for Intersection}]
\label{prop:identity-intersection}
For all $A \subseteq X$: $A \cap X = A$.
\end{proposition}

\begin{remark*}[Fully quantified form]
$\forall A\subseteq X:\ A\cap X = A$.
\end{remark*}

\begin{remark*}[Flash]
\FlashQ{Prove $A\cap X = A$.}
\FlashA{$x\in A\cap X \Leftrightarrow x\in A$ and $x\in X
  \Leftrightarrow x\in A$ and $\top \Leftrightarrow x\in A$.}
\FlashImplies{prop:powerset-boolean-ring}
\end{remark*}

\begin{proposition}[\texorpdfstring{%
  \hyperref[prf:complement-union]{Complement Union is $X$}}{%
  Complement Union is X}]
\label{prop:complement-union}
For all $A \subseteq X$: $A \cup A^c = X$.
\end{proposition}

\begin{remark*}[Fully quantified form]
$\forall A\subseteq X:\ A\cup A^c = X$.
\end{remark*}

\begin{remark*}[Flash]
\FlashQ{Prove $A\cup A^c = X$.}
\FlashA{($\subseteq$): $A\cup A^c\subseteq X$ since both $A,A^c\subseteq X$.
  ($\supseteq$): Let $x\in X$; either $x\in A$ or $x\notin A$, i.e.\ $x\in A^c$.
  So $x\in A\cup A^c$.}
\FlashImplies{prop:powerset-boolean-ring}
\end{remark*}

\begin{proposition}[\texorpdfstring{%
  \hyperref[prf:complement-intersection]{Complement Intersection is $\emptyset$}}{%
  Complement Intersection is Empty}]
\label{prop:complement-intersection}
For all $A \subseteq X$: $A \cap A^c = \emptyset$.
\end{proposition}

\begin{remark*}[Fully quantified form]
$\forall A\subseteq X:\ A\cap A^c = \emptyset$.
\end{remark*}

\begin{remark*}[Flash]
\FlashQ{Prove $A\cap A^c = \emptyset$.}
\FlashA{Suppose $x\in A\cap A^c$; then $x\in A$ and $x\in A^c = X\setminus A$,
  so $x\notin A$. Contradiction. Hence $A\cap A^c = \emptyset$.}
\FlashImplies{prop:powerset-boolean-ring}
\end{remark*}

\begin{proposition}[\texorpdfstring{%
  \hyperref[prf:idempotence-union]{Idempotence of Union}}{Idempotence of Union}]
\label{prop:idempotence-union}
For all $A \subseteq X$: $A \cup A = A$.
\end{proposition}

\begin{remark*}[Fully quantified form]
$\forall A\subseteq X:\ A\cup A = A$.
\end{remark*}

\begin{remark*}[Flash]
\FlashQ{Prove $A\cup A = A$.}
\FlashA{$x\in A\cup A \Leftrightarrow x\in A$ or $x\in A \Leftrightarrow x\in A$.}
\FlashImplies{prop:powerset-boolean-ring}
\end{remark*}

\begin{proposition}[\texorpdfstring{%
  \hyperref[prf:idempotence-intersection]{Idempotence of Intersection}}{%
  Idempotence of Intersection}]
\label{prop:idempotence-intersection}
For all $A \subseteq X$: $A \cap A = A$.
\end{proposition}

\begin{remark*}[Fully quantified form]
$\forall A\subseteq X:\ A\cap A = A$.
\end{remark*}

\begin{remark*}[Flash]
\FlashQ{Prove $A\cap A = A$.}
\FlashA{$x\in A\cap A \Leftrightarrow x\in A$ and $x\in A \Leftrightarrow x\in A$.}
\FlashImplies{prop:powerset-boolean-ring}
\end{remark*}

\begin{proposition}[\texorpdfstring{%
  \hyperref[prf:double-complement]{Double Complement}}{Double Complement}]
\label{prop:double-complement}
For all $A \subseteq X$: $(A^c)^c = A$.
\end{proposition}

\begin{remark*}[Fully quantified form]
$\forall A\subseteq X:\ (A^c)^c = A$.
\end{remark*}

\begin{remark*}[Flash]
\FlashQ{Prove $(A^c)^c = A$.}
\FlashA{$x\in(A^c)^c \Leftrightarrow x\notin A^c \Leftrightarrow
  \neg(x\notin A) \Leftrightarrow x\in A$.}
\FlashImplies{prop:demorgan-cup, prop:demorgan-cap}
\end{remark*}

% ---------------------------------------------------------
% Absorption
% ---------------------------------------------------------

\begin{proposition}[\texorpdfstring{%
  \hyperref[prf:absorption-cup]{Absorption --- Union}}{Absorption --- Union}]
\label{prop:absorption-cup}
For all $A, B \subseteq X$: $A \cup (A \cap B) = A$.
\end{proposition}

\begin{remark*}[Fully quantified form]
$\forall A,B\subseteq X:\ A\cup(A\cap B) = A$.
\end{remark*}

\begin{remark*}[Flash]
\FlashQ{Prove the absorption law $A\cup(A\cap B)=A$.}
\FlashA{($\supseteq$): $A\subseteq A\cup(A\cap B)$ trivially.
  ($\subseteq$): $x\in A\cup(A\cap B)$ means $x\in A$ or ($x\in A$ and $x\in B$).
  In either case $x\in A$.}
\FlashImplies{prop:powerset-boolean-ring}
\end{remark*}

\begin{proposition}[\texorpdfstring{%
  \hyperref[prf:absorption-cap]{Absorption --- Intersection}}{%
  Absorption --- Intersection}]
\label{prop:absorption-cap}
For all $A, B \subseteq X$: $A \cap (A \cup B) = A$.
\end{proposition}

\begin{remark*}[Fully quantified form]
$\forall A,B\subseteq X:\ A\cap(A\cup B) = A$.
\end{remark*}

\begin{remark*}[Flash]
\FlashQ{Prove the absorption law $A\cap(A\cup B)=A$.}
\FlashA{($\subseteq$): $x\in A\cap(A\cup B)$ implies $x\in A$.
  ($\supseteq$): $x\in A$ implies $x\in A\cup B$, hence $x\in A\cap(A\cup B)$.}
\FlashImplies{prop:powerset-boolean-ring}
\end{remark*}

% ---------------------------------------------------------
% Symmetric difference
% ---------------------------------------------------------

\begin{proposition}[\texorpdfstring{%
  \hyperref[prf:symmdiff-comm]{Commutativity of Symmetric Difference}}{%
  Commutativity of Symmetric Difference}]
\label{prop:symmdiff-comm}
For all $A, B \subseteq X$:
$A \mathbin{\triangle} B = B \mathbin{\triangle} A$.
\end{proposition}

\begin{remark*}[Fully quantified form]
$\forall A,B\subseteq X:\ A\triangle B = B\triangle A$.
\end{remark*}

\begin{remark*}[Flash]
\FlashQ{Prove commutativity of symmetric difference.}
\FlashA{$A\triangle B = (A\setminus B)\cup(B\setminus A)
  = (B\setminus A)\cup(A\setminus B) = B\triangle A$,
  using commutativity of union.}
\FlashImplies{prop:powerset-boolean-ring}
\end{remark*}

\begin{proposition}[\texorpdfstring{%
  \hyperref[prf:symmdiff-assoc]{Associativity of Symmetric Difference}}{%
  Associativity of Symmetric Difference}]
\label{prop:symmdiff-assoc}
For all $A, B, C \subseteq X$:
$A \mathbin{\triangle} (B \mathbin{\triangle} C)
= (A \mathbin{\triangle} B) \mathbin{\triangle} C$.
\end{proposition}

\begin{remark*}[Fully quantified form]
$\forall A,B,C\subseteq X:\
  A\triangle(B\triangle C) = (A\triangle B)\triangle C$.
\end{remark*}

\begin{remark*}[Flash]
\FlashQ{What is the proof strategy for associativity of symmetric difference?}
\FlashA{Full element chase: show $x$ belongs to the left side iff it belongs
  to an odd number of the sets $A$, $B$, $C$.  This parity characterisation
  is symmetric in grouping, giving associativity.}
\FlashHint{This is the hardest proof in the Boolean section.  The parity
  argument is the cleanest route: $x\in A\triangle B$ iff $x$ is in exactly
  one of $A$, $B$.}
\FlashImplies{prop:powerset-boolean-ring}
\end{remark*}

\begin{proposition}[\texorpdfstring{%
  \hyperref[prf:symmdiff-self]{Symmetric Difference with Self}}{%
  Symmetric Difference with Self}]
\label{prop:symmdiff-self}
For all $A \subseteq X$: $A \mathbin{\triangle} A = \emptyset$.
\end{proposition}

\begin{remark*}[Fully quantified form]
$\forall A\subseteq X:\ A\triangle A = \emptyset$.
\end{remark*}

\begin{remark*}[Flash]
\FlashQ{Prove $A\triangle A = \emptyset$.}
\FlashA{$A\triangle A = (A\setminus A)\cup(A\setminus A) = \emptyset\cup\emptyset = \emptyset$.}
\FlashImplies{prop:powerset-boolean-ring}
\end{remark*}

\begin{proposition}[\texorpdfstring{%
  \hyperref[prf:symmdiff-empty]{Symmetric Difference with Empty Set}}{%
  Symmetric Difference with Empty Set}]
\label{prop:symmdiff-empty}
For all $A \subseteq X$: $A \mathbin{\triangle} \emptyset = A$.
\end{proposition}

\begin{remark*}[Fully quantified form]
$\forall A\subseteq X:\ A\triangle\emptyset = A$.
\end{remark*}

\begin{remark*}[Flash]
\FlashQ{Prove $A\triangle\emptyset = A$.}
\FlashA{$A\triangle\emptyset
  = (A\setminus\emptyset)\cup(\emptyset\setminus A)
  = A\cup\emptyset = A$.}
\FlashImplies{prop:powerset-boolean-ring}
\end{remark*}

\begin{proposition}[\texorpdfstring{%
  \hyperref[prf:cap-dist-symmdiff]{%
    Distributivity of $\cap$ over $\triangle$}}{%
  Distributivity of Intersection over Symmetric Difference}]
\label{prop:cap-dist-symmdiff}
For all $A, B, C \subseteq X$:
$A \cap (B \mathbin{\triangle} C)
= (A \cap B) \mathbin{\triangle} (A \cap C)$.
\end{proposition}

\begin{remark*}[Fully quantified form]
$\forall A,B,C\subseteq X:\
  A\cap(B\triangle C) = (A\cap B)\triangle(A\cap C)$.
\end{remark*}

\begin{remark*}[Flash]
\FlashQ{State and outline the proof that $\cap$ distributes over $\triangle$.}
\FlashA{$A\cap(B\triangle C) = (A\cap B)\triangle(A\cap C)$.
  Element chase: $x\in A\cap(B\triangle C)$ iff $x\in A$ and $x$ is in exactly
  one of $B$, $C$, iff $x$ is in exactly one of $A\cap B$, $A\cap C$,
  iff $x\in(A\cap B)\triangle(A\cap C)$.}
\FlashHint{This is the ring distributivity law: in the Boolean ring
  $(\mathcal{P}(X),\triangle,\cap)$, $\cap$ is multiplication distributing
  over $\triangle$ as addition.}
\FlashImplies{prop:powerset-boolean-ring}
\end{remark*}

% ---------------------------------------------------------
% Boolean ring capstone
% ---------------------------------------------------------

\begin{tcolorbox}[colback=thmbox, colframe=thmborder, arc=2pt,
  left=6pt, right=6pt, top=4pt, bottom=4pt,
  title={\small\textbf{Theorem (Power Set Boolean Ring)}},
  fonttitle=\small\bfseries]
\label{thm:powerset-boolean-ring}
Let $X$ be a non-empty set.  Then
$(\mathcal{P}(X),\mathbin{\triangle},\cap)$ is a commutative ring
with unity in which every element is idempotent (a Boolean ring).
The additive identity is $\emptyset$ and the multiplicative identity
is $X$.
\end{tcolorbox}

\begin{proposition}[\texorpdfstring{%
  \hyperref[prf:powerset-boolean-ring]{%
    $(\mathcal{P}(X),\triangle,\cap)$ is a Boolean Ring}}{%
  Power Set is a Boolean Ring}]
\label{prop:powerset-boolean-ring}
$(\mathcal{P}(X), \mathbin{\triangle}, \cap)$ is a Boolean ring.
\end{proposition}

\begin{remark*}[Fully quantified form]
$\forall$ non-empty set $X$:
$(\mathcal{P}(X),\triangle,\cap)$ satisfies all ring axioms,
is commutative, has unity $X$, and satisfies $A\cap A = A$
for all $A\in\mathcal{P}(X)$.
\end{remark*}

\begin{remark*}[Flash]
\FlashQ{What ring structure does $\mathcal{P}(X)$ carry, and what are the
  additive and multiplicative identities?}
\FlashA{$(\mathcal{P}(X),\triangle,\cap)$ is a Boolean ring.
  Additive identity: $\emptyset$ (since $A\triangle\emptyset = A$).
  Multiplicative identity: $X$ (since $A\cap X = A$).
  Every element satisfies $A\cap A = A$ (idempotence).}
\FlashHint{The ring axioms are assembled from the Boolean law propositions
  proved earlier in this section.  No new element chasing is required ---
  cite the named propositions.}
\FlashImplies{def:ring-of-sets, def:algebra-of-sets}
\end{remark*}
